% ============================================================================
% main.tex — Masterarbeit im Layout der Universität Oldenburg (MA-BA-Vorlage)
%
% Struktur und Optionen folgen der Vorlage da.tex (scrreprt, BCOR12mm,
% twoside, openright, Titelseite.inc, Danksagung, Erklärung, Zusammenfassung,
% Verzeichnisse). Anpassungen gegenüber der Original-Vorlage:
%   - scrpage2 -> scrlayer-scrpage (scrpage2 ist in aktuellen TeX-Live-
%     Distributionen nicht mehr enthalten; identische Kopf-/Fußzeilen)
%   - bibgerm/geralpha + BibTeX -> formatiertes Literaturverzeichnis als
%     Kapitel im Text (die Arbeit führt ihre Referenzen im APA-Stil)
%   - veraltete Pakete (epsfig, ae/aecompl, subfigure) entfernt
%   - Pakete für die aus R Markdown erzeugten Tabellen/Abbildungen ergänzt
%     (booktabs, longtable, threeparttable, xcolor[table], float, ...)
%   - babel: Hauptsprache Englisch (Text), Deutsch für Formalseiten
%
% Kompilieren AUS DEM REPO-ROOT (wegen der Abbildungspfade):
%   pdflatex -output-directory=thesis thesis/main.tex   (2x ausführen)
% ============================================================================
\documentclass[a4paper,titlepage,liststotoc,final,twoside,openright,cleardoubleempty,BCOR12mm]{scrreprt}

\usepackage[utf8]{inputenc}
\usepackage[T1]{fontenc}
\usepackage[ngerman,english]{babel}   % Englisch = Hauptsprache des Texts
\usepackage{amsmath}
\usepackage{amssymb}
\usepackage[table]{xcolor}
\usepackage{graphicx}
\graphicspath{{./}{thesis/}}
\usepackage{tabularx}
\usepackage{rotating}
\usepackage{lscape}

% Tabellen/Abbildungen aus R Markdown (kableExtra/knitr)
\usepackage{booktabs}
\usepackage{longtable}
\usepackage{array}
\usepackage{multirow}
\usepackage{threeparttable}
\usepackage{caption}
\captionsetup[table]{skip=6pt}
\usepackage{float}
\floatplacement{figure}{H}
\usepackage{placeins}
\renewcommand{\arraystretch}{1.15}

% Pandoc-Hilfsmakros (vom R-Markdown-Body benötigt)
\providecommand{\tightlist}{\setlength{\itemsep}{0pt}\setlength{\parskip}{0pt}}
\makeatletter
\newsavebox\pandoc@box
\newcommand*\pandocbounded[1]{% Bilder auf Seitengröße begrenzen
  \sbox\pandoc@box{#1}%
  \Gscale@div\@tempa{\textheight}{\dimexpr\ht\pandoc@box+\dp\pandoc@box\relax}%
  \Gscale@div\@tempb{\linewidth}{\wd\pandoc@box}%
  \ifdim\@tempb\p@<\@tempa\p@\let\@tempa\@tempb\fi
  \ifdim\@tempa\p@<\p@\scalebox{\@tempa}{\usebox\pandoc@box}%
  \else\usebox{\pandoc@box}%
  \fi
}
\makeatother
\newcommand{\real}[1]{#1}

\usepackage{hyperref}
\hypersetup{pdftitle={Demand Estimation with Unstructured Product Data},
            pdfauthor={Burak Uzun}, hidelinks}

% Vorlage – Obere Linie, Seitenzahlen (scrlayer-scrpage statt scrpage2)
\usepackage[automark,headsepline=.4pt]{scrlayer-scrpage}
\pagestyle{scrheadings}
\clearpairofpagestyles
\automark[section]{chapter}
\ihead{\leftmark}
\ohead{\rightmark}
\ofoot*{\pagemark}

% Arbeitsbeschreibung
\title{Demand Estimation with Unstructured Product Data: Evidence from Amazon's Toothpaste Market}
\author{Burak Uzun}
\date{\today}

\begin{document}

% Fuers Inhaltsvz gibts andere Seitenzahlen
\pagenumbering{Roman}

% Titelseite erstellen
\begin{titlepage}
 \input{thesis/Titelseite.inc}
\end{titlepage}

% Leere Seite
\newpage
\thispagestyle{empty}
\quad

% Danksagung & Eigenstaendigkeitserklaerung
\newpage
\thispagestyle{empty}
\quad
\begin{otherlanguage}{ngerman}
\section*{Danksagung}
% TODO: Danksagung einfuegen (optional)
[Danksagung hier einf\"ugen.]

\vspace{6.8cm}
\begin{center}
\begin{Large}
\textbf{Erkl\"arung}
\end{Large}
\end{center}

\noindent
Ich erkl\"are, dass ich diese Arbeit selbstst\"andig verfasst und keine anderen als die angegebenen
Quellen und Hilfsmittel verwendet habe.

\vspace{3cm}
\noindent
Burak Uzun\\
% TODO: Matrikelnummer eintragen
Matrikelnummer [XXXXXXXX]\\
Oldenburg, den [TT]. [Monat] 2026
\end{otherlanguage}

% Zusammenfassung erstellen (jeweils auf Englisch und Deutsch)
\begin{abstract}
\section*{Abstract}
Merger simulation hinges on correctly measuring substitution patterns, yet the
nesting structures used in practice are typically built from coarse labels such
as brand identity. This thesis asks whether product representations extracted
from unstructured listing content -- images and text encoded with CLIP -- define
substitution patterns better than brand labels do. Using an ASIN-by-quarter
panel of the Amazon US toothpaste market around Colgate-Palmolive's 2020
acquisition of the natural-brand hello, three nested-logit demand models are
estimated via the Berry (1994) linearisation that differ only in how nests are
assigned: no nesting, brand nesting, and nesting by k-means clusters of CLIP
image-and-text embeddings. CLIP-based nesting fits demand best and cuts the
share of economically implausible negative implied marginal costs from 8.6\% to
1.7\%. Because CLIP places hello and Colgate in different clusters, the implied
diversion between the merging parties is small: the simulated merger price
effect for hello falls from +1.23\% under plain logit to +0.40\% under CLIP
nesting. The results suggest embedding-based nests are a practical, replicable
tool for antitrust analysis when conventional product characteristics are
unavailable.

\begin{otherlanguage}{ngerman}
\section*{Zusammenfassung}
Fusionssimulationen h\"angen entscheidend davon ab, Substitutionsmuster korrekt
zu messen; die in der Praxis verwendeten Nest-Strukturen beruhen jedoch meist
auf groben Merkmalen wie der Markenzugeh\"origkeit. Diese Arbeit untersucht, ob
Produktrepr\"asentationen aus unstrukturierten Angebotsinhalten -- mit CLIP
kodierte Bilder und Texte -- Substitutionsmuster besser abbilden als
Markenlabels. Auf Basis eines ASIN-Quartals-Panels des US-amerikanischen
Amazon-Zahnpastamarktes rund um die \"Ubernahme der Naturkosmetikmarke hello
durch Colgate-Palmolive im Jahr 2020 werden drei Nested-Logit-Nachfragemodelle
mittels der Berry-(1994)-Linearisierung gesch\"atzt, die sich ausschlie{\ss}lich
in der Nest-Zuordnung unterscheiden: ohne Nester, Marken-Nester und Nester aus
k-Means-Clustern der CLIP-Einbettungen. Die CLIP-basierte Nest-Struktur weist
die beste Anpassung auf und senkt den Anteil \"okonomisch unplausibler negativer
implizierter Grenzkosten von 8{,}6\,\% auf 1{,}7\,\%. Da CLIP hello und Colgate
unterschiedlichen Clustern zuordnet, ist die implizierte Diversion zwischen den
Fusionsparteien gering: Der simulierte Preiseffekt der Fusion f\"ur hello sinkt
von +1{,}23\,\% (einfaches Logit) auf +0{,}40\,\% (CLIP-Nester). Die Ergebnisse
legen nahe, dass einbettungsbasierte Nester ein praktikables, reproduzierbares
Instrument der Wettbewerbsanalyse sind, wenn klassische Produkteigenschaften
nicht verf\"ugbar sind.
\end{otherlanguage}
\end{abstract}

% Inhaltsverzeichnis erstellen
\tableofcontents
% Abbildungsverzeichnis
\listoffigures
% Tabellenverzeichnis
\listoftables

\cleardoublepage
\pagenumbering{arabic} % Seitenzahlen umstellen

% ---- Inhalt (aus blp_thesis.Rmd generiert; nicht von Hand editieren) ----



\begin{center}\rule{0.5\linewidth}{0.5pt}\end{center}

\chapter{Introduction}\label{introduction}

Antitrust authorities increasingly confront mergers in digital markets
-- characteristically, an established incumbent acquiring a
fast-growing, digitally native brand -- in which product differentiation
is not described by standardised attributes. In a car merger the analyst
observes horsepower and fuel economy; in a laptop merger, RAM and screen
size. In an online consumer-goods merger, the economically relevant
differentiation -- a ``natural'' positioning, sensitivity relief, visual
branding, a kid-friendly format -- lives in unstructured listing
content: product images, titles, and free-form benefit descriptions. The
workhorse tools of merger simulation, however, were built for the
structured world. Nested-logit models (McFadden 1978) encode
substitution through nest assignments constructed from coarse observable
labels, most commonly brand identity or hand-picked segments, and when
those labels misrepresent how consumers perceive products, the resulting
diversion ratios, markups, and simulated price effects inherit the
error.

A second, equally practical obstacle compounds the first. Estimating
price sensitivity credibly requires exogenous cost-side variation, and
high-frequency online retail data rarely contain any: input-cost series
vary only over time and are absorbed by time fixed effects, while
Hausman-style rival-price instruments (Hausman 1996) fail when demand
shocks are correlated across markets, as they plainly are on a single
national platform. The symptom of proceeding without valid instruments
is well known to practitioners: demand estimated far too inelastic,
implying Bertrand markups so large that the marginal costs recovered
from the first-order conditions turn \emph{negative} for a substantial
share of products. This negative-cost rate is not a cosmetic nuisance --
it is a direct, observable measure of how badly a demand system
describes competition, and it is the diagnostic this paper puts at the
centre of the analysis.

This paper offers a pragmatic pipeline for exactly that environment,
with one methodological contribution at its core. The pipeline maps
unstructured product listings into a continuous differentiation space
using a pre-trained deep learning model -- CLIP image and text
embeddings (Radford et al.~2021), reduced to five joint principal
components -- and sidesteps the missing-instrument problem by
calibrating the price coefficient to an established industry elasticity,
the meta-analytic mean own-price elasticity of \(-2.62\) for consumer
packaged goods (Bijmolt, van Heerde, and Pieters 2005). The contribution
lies in what the embeddings are used \emph{for}: nests are assigned by
k-means clustering in the embedding space rather than by brand. Holding
everything else fixed, we compare three Berry (1994) nested-logit models
that differ only in their nest assignment -- no nesting, brand nesting,
and CLIP-cluster nesting -- and let the negative-cost diagnostic
adjudicate between them.

The empirical laboratory is Colgate-Palmolive's acquisition of hello
Products LLC, a fast-growing premium ``naturally friendly'' oral-care
brand. The deal was announced on 23 January 2020 and completed on 31
January 2020 for \$351 million in cash, giving a clean event date
(2020Q1) inside the sample. Using an ASIN-by-quarter panel of the Amazon
US toothpaste market (2018Q1--2022Q3, 159 products), we estimate three
demand models via the Berry (1994) nested-logit linearisation that are
identical in every respect except the nest assignment: M1 imposes no
nesting (plain logit), M2 nests by brand, and M3 nests by K=5 CLIP
clusters. With the price coefficient fixed by calibration, the remaining
parameters are estimated by concentrated OLS with quarter fixed effects
and the nesting parameter \(\rho\) is profiled by grid search. Marginal
costs are recovered from the Bertrand first-order conditions, and the
merger is simulated by recomputing Nash equilibrium prices under pre-
and post-merger ownership at fixed pre-merger costs.

Three results stand out. First, CLIP-based nesting fits demand best: the
profiled nesting parameter is \(\rho^{*}=0.67\) and the residual sum of
squares falls by 8.1\% relative to plain logit, against 3.1\% for brand
nesting. Second -- and this is the central finding -- CLIP nesting
drives the share of observations with \emph{negative implied marginal
costs} from 8.6\% under both plain logit and brand nesting to 1.7\%.
Negative recovered costs are the canonical symptom of a demand system
that understates substitutability and therefore overstates markups;
their near-elimination is direct evidence that embedding-based nests
capture the competitive structure that brand labels miss. Third, the
economics of the merger change materially: CLIP places hello and Colgate
in different clusters, implying low diversion between the merging
parties, and the simulated price effect for hello falls from +1.23\%
(plain logit) to +0.40\% (CLIP nesting), with Colgate's own effect
essentially zero.

Beyond the specific merger, the exercise is designed as a transparent,
fully reproducible template for antitrust practitioners: every step --
from embedding construction through the \(\rho\) grid search to the Nash
counterfactuals -- runs from a single document, and the Nash solver is
validated against PyBLP's counterfactual routines (Conlon and Gortmaker
2020) to numerical precision. The remainder of the paper proceeds as
follows. Section 2 relates the approach to the literature. Section 3
describes the industry setting, the data, and the embedding
construction. Section 4 lays out the empirical framework. Section 5
presents the results, and Section 6 discusses limitations. Section 7
concludes.

\begin{center}\rule{0.5\linewidth}{0.5pt}\end{center}

\chapter{Related Literature}\label{related-literature}

The demand framework descends from the discrete-choice literature on
differentiated products. Berry (1994) shows that logit and nested-logit
demand can be estimated by linear methods after inverting market shares,
and that the same first-order conditions used to rationalise observed
prices can be inverted to recover marginal costs under Bertrand
competition. Berry, Levinsohn, and Pakes (1995) generalise this
framework to random-coefficients demand; Grigolon and Verboven (2014)
show that a well-specified nested logit often approximates
random-coefficients substitution patterns closely, which makes the
choice of nests -- the object studied here -- the binding constraint in
practice. Nevo (2000a) provides the canonical practitioner's treatment
of this class of models, and Nevo (2001) demonstrates the full
structural program -- demand estimates, recovered margins, and
competitive counterfactuals -- in the ready-to-eat cereal industry, a
consumer packaged goods setting closely analogous to toothpaste. Merger
simulation itself follows Werden and Froeb (1994) and Nevo (2000b):
recover marginal costs from the pre-merger equilibrium, change the
ownership structure, and re-solve for Bertrand-Nash prices -- computed
here with the damped fixed-point iteration analysed by Morrow and
Skerlos (2011). Conlon and Gortmaker (2020) codify current best practice
and provide the PyBLP routines against which this paper's equilibrium
computations are validated. Draganska and Jain (2006) estimate
structural demand for differentiated CPG product lines and report
own-price elasticities of the magnitude used for calibration here.

Within this tradition, the paper's price coefficient is calibrated
rather than instrumented. The target elasticity of \(-2.62\) comes from
Bijmolt, van Heerde, and Pieters (2005), who meta-analyse 1,851
own-price elasticity estimates from 81 studies and report \(-2.62\) as
the mean. Calibration to an external elasticity is standard in applied
antitrust simulation when credible cost-side instruments are unavailable
(the logit merger calibrations of Werden and Froeb 1994 are the
template), and it has a transparency advantage in a practitioner
context: the price-sensitivity assumption is stated in one number rather
than buried in a first stage.

The use of unstructured data is most closely related to Compiani,
Morozov, and Seiler (2026), who extract embeddings from product images
and text with pre-trained deep learning models and incorporate them into
a mixed logit model, showing across 40 Amazon categories that
unstructured data identify close substitutes that attribute-based models
miss. This paper shares their premise -- embeddings encode the
differentiation consumers care about -- but differs in implementation
and aim. Where Compiani, Morozov, and Seiler (2026) enter embeddings as
continuous random-coefficient characteristics in a mixed logit, we use
them to \emph{define the nesting structure} of a Berry (1994) nested
logit. The cost of this simplification is a coarser substitution
pattern; the benefit is that the entire pipeline remains linear,
transparent, and estimable in seconds, which matters for the
merger-review setting where speed and auditability are first-order
concerns. The embedding technology itself, CLIP, is due to Radford et
al.~(2021); its contrastive image-text training makes image and text
representations directly comparable, which is what allows a single joint
product space to be built from both modalities.

The contribution, then, sits at the intersection: a minimal, replicable
bridge between modern representation learning and the nested-logit
merger simulation toolkit that antitrust practitioners already use.

\begin{center}\rule{0.5\linewidth}{0.5pt}\end{center}

\chapter{Industry Background and
Data}\label{industry-background-and-data}

\section{The hello--Colgate
Acquisition}\label{the-hellocolgate-acquisition}

hello Products LLC, founded in 2012, positioned itself as a ``naturally
friendly'' oral-care brand -- fluoride-optional formulas, charcoal and
coconut-oil lines, playful branding -- with strong appeal among younger
consumers. Colgate-Palmolive announced its agreement to acquire hello on
23 January 2020 and completed the transaction on 31 January 2020 for
\$351 million in cash. The acquisition is an attractive empirical
laboratory for three reasons: the event date falls cleanly at the start
of 2020Q1; the acquirer is the category's largest incumbent while the
target occupies a differentiated ``natural'' niche, so the price effect
hinges on exactly the substitution question this paper studies; and both
parties sell extensively on Amazon, where listing content is observable.

\section{Data Sources}\label{data-sources}

The analysis combines two data sources. Household-level purchase
histories come from the Open E-commerce 1.0 dataset (Berke et al.~2024),
a crowdsourced panel of 5,027 US Amazon consumers with more than 1.8
million transaction records, distributed through Harvard Dataverse. The
dataset has been validated by its authors against public Amazon sales
data: aggregate expenditure in the panel correlates with published
Amazon revenues at Pearson \(r = 0.978\) (\(p < 0.001\); Berke et
al.~2024). The toothpaste purchase records extracted for this study span
January 2018 to March 2023 and include order timestamps, prices, and
household identifiers; household demographics are available in the
source data but enter the analysis only through the count of unique
purchasing households per cell. Product-level metadata -- titles,
benefit descriptions, and image URLs for every purchased ASIN -- were
retrieved programmatically through the Keepa API (a subscription-based
Amazon data service) and parsed with a rule-based extraction script
(brand matching against a curated brand list, keyword-based benefit
extraction, and unit-aware package-size parsing). Prices are deflated to
real terms using the CPI before aggregation; the panel's quarterly price
for each ASIN is the mean deflated transaction price across that
quarter's purchases.

\section{Sample Construction}\label{sample-construction}

The estimation sample is built in three steps, each dictated by a
practical constraint:

\begin{enumerate}
\def\labelenumi{\arabic{enumi}.}
\item
  \textbf{Choice-set restriction (167 ASINs).} The universe of
  toothpaste ASINs is restricted to products with at least five recorded
  purchases over the sample window. An unrestricted choice set with
  thousands of rarely purchased alternatives proved computationally
  infeasible and produced singular design matrices; the five-purchase
  floor keeps every product with meaningful revealed-preference
  information while discarding noise listings.
\item
  \textbf{Brand grouping (12 groups).} Brands accounting for more than
  2\% of purchases are treated as their own groups -- Colgate, Crest,
  SENSODYNE PRONAMEL, Sensodyne, Tom's of Maine, Arm \& Hammer, hello,
  Parodontax, APAGARD, JASON, and Orajel -- and all remaining small
  brands are collapsed into an ``Other'' category. Together the eleven
  named groups account for 87\% of purchases.
\item
  \textbf{Panel filters (159 ASINs, 1,164 observations).} Purchases are
  aggregated to ASIN x quarter cells. The estimation window ends in
  2022Q3 (later quarters are incompletely covered), and ASINs must
  appear in at least three quarters to enter the sample, eliminating
  products whose near-zero shares would be dominated by sampling noise.
  The resulting panel covers 2018Q1--2022Q3: 19 quarters, 159 ASINs, and
  1,164 observations.
\end{enumerate}

Market shares are constructed from revealed-preference purchase counts.
Market size \(M_t\) is set to \((1 + \lambda) \cdot \sum_j q_{jt}\) with
\(\lambda = 13.93\), implying an outside-good share of approximately
93\%.

Separately from these sample filters, the joint embedding space of
Section 3.4 is \emph{fitted} on the 155 of 167 ASINs that hold both a
valid image and a valid text embedding; the 12 ASINs with text-only
embeddings are projected into that space and remain in the estimation
sample throughout. Table 1 reports summary statistics for the estimation
panel.

\begin{table}[H]
\centering
\caption{\label{tab:tab_panel_summary}Panel Summary Statistics (Estimation Sample)}
\centering
\fontsize{9}{11}\selectfont
\begin{threeparttable}
\begin{tabular}[t]{lrrrrr}
\toprule
Variable & Mean & Median & SD & Min & Max\\
\midrule
Price (USD) & 9.20 & 8.47 & 5.34 & 1.35 & 41.75\\
Quarterly units sold & 2.17 & 1.00 & 1.96 & 1.00 & 18.00\\
Market share (x1,000) & 1.093 & 0.753 & 1.036 & 0.335 & 9.802\\
Package size (g) & 320 & 259 & 223 & 43 & 1179\\
Purchasing households per cell & 2.06 & 1.00 & 1.85 & 1.00 & 17.00\\
\bottomrule
\end{tabular}
\begin{tablenotes}
\item \textit{Notes: } 
\item 1,164 ASIN x quarter observations; 159 ASINs; 19 quarters (2018Q1--2022Q3); 12 brand groups. Price is the quarterly mean deflated transaction price per ASIN. Market share uses market size M\_t = (1 + 13.93) x total quarterly purchases (outside-good share approx. 93\%).
\end{tablenotes}
\end{threeparttable}
\end{table}

\section{From Listings to Nests: Embeddings, Joint PCA, and
Clustering}\label{from-listings-to-nests-embeddings-joint-pca-and-clustering}

Each product \(j\) is represented by its listing image and the
concatenation of its title and benefit description. CLIP (Contrastive
Language-Image Pre-training; Radford et al.~2021) provides an image
encoder \(f_I\) and a text encoder \(f_T\), trained contrastively so
that matching image--text pairs lie close together in a shared
representation space. Both encoders (\texttt{clip-vit-base-patch32}, a
Vision Transformer; Dosovitskiy et al.~2021) map into
\(\mathbb{R}^{512}\), and outputs are L2-normalised:

\[\mathbf{e}^{I}_{j} = \frac{f_I(\text{image}_j)}{\lVert f_I(\text{image}_j) \rVert}, \qquad
\mathbf{e}^{T}_{j} = \frac{f_T(\text{text}_j)}{\lVert f_T(\text{text}_j) \rVert}.\]

CLIP's usefulness for this application comes from how it was trained.
The model consists of two separate networks: a text Transformer, and --
in the variant used here -- the ViT-B/32 Vision Transformer, which cuts
each image into 32 x 32-pixel patches and processes the patch sequence
the way a language model processes tokens (Radford et al.~2021,
sec.~2.4; Dosovitskiy et al.~2021). Both networks project into the same
512-dimensional space, and training proceeds contrastively on 400
million (image, caption) pairs collected from the web: within each
training batch of \(N\) pairs, the model is rewarded for maximising the
cosine similarity of the \(N\) correct image--caption pairings while
minimising the similarity of all mismatched pairings (Radford et
al.~2021, sec.~2; their Figure 1 summarises the scheme). Figure 1
sketches this objective and its role in the empirical pipeline.

Two properties of the contrastive objective matter here. First, because
web captions describe products in ordinary marketing language
(``charcoal whitening toothpaste'', ``fluoride-free kids toothpaste''),
the learned space encodes exactly the attribute vocabulary that
differentiates consumer goods, with no category-specific training
required. Second, because images and text are embedded in one shared
space, the two modalities can be combined into a single product
representation -- which is what the joint PCA below exploits.

\begin{figure}
\centering
\includegraphics{thesis_body_files/figure-latex/fig_clip_scheme-1.pdf}
\caption{How CLIP embeddings are produced and used. Left: the
contrastive training objective (after Radford et al.~2021, their Figure
1). An image encoder and a text encoder map a batch of image-caption
pairs into a shared space; training pushes the correct pairings (dark
diagonal cells) toward high cosine similarity and all mismatched
pairings (light cells) apart. Right: the empirical pipeline of this
paper, from raw listing content to the nests of model M3.}
\end{figure}

Each modality's embeddings are first reduced to five modality-specific
principal components (Hotelling 1933), with each PCA fitted only on the
products holding a valid embedding for that modality. This yields a
ten-dimensional characteristics vector per product,
\(\mathbf{x}_j = (x^{I}_{j1},\dots,x^{I}_{j5},\,x^{T}_{j1},\dots,x^{T}_{j5})\).

Because these ten columns come from two separate PCA fits, their scales
are not comparable: within each modality the first component
mechanically carries the largest variance, and the image and text blocks
have different overall variance levels. The joint PCA is therefore
computed on the correlation scale. Each column \(k\) is z-scored using
its mean \(\mu_k\) and standard deviation \(\sigma_k\) over the
\(N = 155\) products with both modalities,

\[z_{jk} = \frac{x_{jk} - \mu_k}{\sigma_k},\]

and the rotation matrix \(\mathbf{W}\) is obtained from the PCA of the
standardised matrix \(\mathbf{Z}\). The first five joint components,
\(\mathbf{pc}_j = \mathbf{z}_j' \mathbf{W}_{[\cdot,1:5]}\), define the
differentiation space used throughout the paper. Standardisation is not
cosmetic: without it, the joint solution would be dominated by whichever
input columns happen to carry the largest raw variance -- the two
first-stage PC1s -- rather than by genuine cross-modality structure. The
fitted parameters \((\mu_k, \sigma_k, \mathbf{W})\) are saved with the
replication materials, which makes the projection of any product into
the joint space exactly reproducible.

A useful diagnostic for whether the joint space actually blends the two
modalities is the modality share of each component: the fraction of its
squared rotation loadings attributable to image inputs,

\[\text{share}^{I}_{m} = \frac{\sum_{k \in I} w_{km}^2}{\sum_{k} w_{km}^2}, \qquad m = 1, \dots, 5.\]

As Figure 2 shows, every retained joint component draws between 48.6\%
and 51.4\% of its loading mass from each modality -- the joint space is
genuinely multimodal rather than a relabelled image space or text space.

Finally, nests are formed by k-means clustering (Lloyd 1982) on the five
joint components: each product is assigned to one of \(K = 5\) clusters,

\[g_j = \arg\min_{g \in \{1,\dots,K\}} \lVert \mathbf{pc}_j - \mathbf{c}_g \rVert^2,\]

with centroids \(\mathbf{c}_g\) from 50 random restarts (seed 42), and
cluster membership is the nest assignment in M3. The five clusters
explain 71.2\% of joint-PC variance. Appendix A reports cluster-quality
diagnostics -- within-cluster sum of squares and average silhouette
width (Rousseeuw 1987) -- together with all downstream results for
\(K = 2, \dots, 10\).

\begin{figure}
\centering
\includegraphics{thesis_body_files/figure-latex/fig_pca-1.pdf}
\caption{Anatomy of the joint embedding space. Left: variance explained
by each joint principal component of the standardised ten-column input
(five image PCs and five text PCs), fitted on the 155 ASINs with both
modalities; the first five components (blue) are retained. Right:
modality composition of each retained component, measured as the share
of squared rotation loadings attributable to image versus text inputs.
The dashed line marks an even 50/50 split; every component lies within
1.4 percentage points of it, confirming the joint space blends both
modalities rather than being dominated by either.}
\end{figure}

Embeddings are static over the sample window. Established Amazon
listings change their images and titles only marginally over 2018--2022,
so re-extracting embeddings per quarter would produce near-identical
vectors; the identifying variation is \textbf{cross-sectional} --- ASINs
within the same brand differ genuinely in packaging and benefit claims
--- and this is exactly the variation recovered by moving from
brand-level to ASIN-level data. Time-varying demand conditions are
absorbed by quarter fixed effects.

Twelve ASINs lack an image embedding (image retrieval failed at
extraction time) but have valid text embeddings. Rather than dropping
them or imputing brand means, they are projected into the joint PC space
from their standardised text components, with image dimensions set to
the cross-sectional mean -- zero in standardised space, so that
\(\mathbf{pc}_j = (\mathbf{0},\, \mathbf{z}^{T}_j)'\,\mathbf{W}_{[\cdot,1:5]}\)
using the same saved \((\mu_k, \sigma_k, \mathbf{W})\). These ASINs span
Crest (4), Sensodyne (2), Tom's of Maine (2), hello (2), and Other (2),
and all clear the three-quarter inclusion filter. Dropping them would
distort brand-level market shares and remove hello observations inside
the merger simulation window; brand-mean imputation was rejected because
it would reintroduce precisely the aggregation that the ASIN-level
design is meant to avoid.

\section{The Estimated Product Space}\label{the-estimated-product-space}

Table 2 reports summary statistics by CLIP cluster. The clusters recover
economically meaningful segments without using any sales or price
information: a premium sensitivity segment (C1, dominated by the GSK
brands), the two mainstream incumbents in their own clusters (C2
Colgate, C3 Crest), Tom's of Maine (C4), and a natural/specialty segment
(C5) containing hello. Price levels and package sizes differ sharply
across clusters, consistent with the clusters capturing vertical and
horizontal differentiation simultaneously.

\begin{table}[!h]
\centering
\caption{\label{tab:table1}Summary Statistics by CLIP Cluster}
\centering
\resizebox{\ifdim\width>\linewidth\linewidth\else\width\fi}{!}{
\fontsize{8}{10}\selectfont
\begin{threeparttable}
\begin{tabular}[t]{llrrrrrrrrrr}
\toprule
\multicolumn{3}{c}{ } & \multicolumn{3}{c}{Price (USD)} & \multicolumn{3}{c}{Mkt. Share (x1000)} & \multicolumn{3}{c}{Pkg. Size (g)} \\
\cmidrule(l{3pt}r{3pt}){4-6} \cmidrule(l{3pt}r{3pt}){7-9} \cmidrule(l{3pt}r{3pt}){10-12}
Cluster & Dom. Brand & $N$ ASINs & $\bar{p}$ & $\tilde{p}$ & $\sigma_p$ & $\bar{s}\times10^3$ & $\tilde{s}\times10^3$ & $\sigma_s\times10^3$ & $\bar{w}$ & $\tilde{w}$ & $\sigma_w$\\
\midrule
C1 & Sensodyne & 29 & 10.26 & 9.63 & 5.58 & 1.068 & 0.770 & 0.848 & 249 & 249 & 145\\
C2 & Colgate & 33 & 8.00 & 8.39 & 3.65 & 1.303 & 0.817 & 1.259 & 537 & 399 & 296\\
C3 & Crest & 29 & 8.72 & 6.88 & 4.18 & 1.267 & 0.757 & 1.282 & 375 & 431 & 128\\
C4 & Tom's of Maine & 13 & 9.85 & 9.76 & 2.96 & 1.170 & 0.770 & 0.972 & 374 & 370 & 94\\
C5 & Other (hello) & 55 & 9.46 & 7.86 & 6.84 & 0.839 & 0.670 & 0.736 & 171 & 141 & 105\\
\bottomrule
\end{tabular}
\begin{tablenotes}
\item \textit{Notes: } 
\item Panel: 1,164 ASIN x quarter observations, 19 quarters (2018Q1--2022Q3). Mean/Median/SD of price in USD, market share scaled x1000, package weight in grams. Dominant brand = modal brand by ASIN count within cluster.
\end{tablenotes}
\end{threeparttable}}
\end{table}

\vspace{1em}

Beyond price and share statistics, the clusters are directly
interpretable from the listing content itself. Table 3 reports, for each
cluster, the listing-vocabulary terms that are most over-represented
relative to the whole category (by lift: the share of the cluster's
ASINs whose title or benefit text contains the term, divided by the
category-wide share). The vocabulary confirms that the clusters track
economically meaningful positioning rather than arbitrary embedding
geometry.

\begin{table}[H]
\centering
\caption{\label{tab:tab_vocab}Distinctive Listing Vocabulary by CLIP Cluster}
\centering
\fontsize{9}{11}\selectfont
\begin{threeparttable}
\begin{tabular}[t]{llr>{\raggedright\arraybackslash}p{22em}}
\toprule
Cluster & Dominant brand & N ASINs & Most over-represented listing terms (by lift)\\
\midrule
C1 & Sensodyne & 29 & sensodyne, parodontax, extra, bleeding, pronamel, reharden\\
C2 & Colgate & 33 & optic, hydrogen, max, benefits, extreme, frosty\\
C3 & Crest & 29 & scope, pro, radiant, brilliance, eraser, glamorous\\
C4 & Tom's of Maine & 13 & maine, tom's, children's, silly, animals, not\\
C5 & Other & 55 & sls, activated, mouth, orajel, remineralizing, anti\\
\bottomrule
\end{tabular}
\begin{tablenotes}
\item \textit{Notes: } 
\item Terms from ASIN titles and benefit descriptions (lower-cased, unit and category stopwords removed). Lift = share of the cluster's ASINs containing the term divided by the category-wide share; terms shown appear in at least two of the cluster's ASINs and have lift above 1.5, ranked by lift. Cluster assignment identical to the estimation (K = 5, seed 42).
\end{tablenotes}
\end{threeparttable}
\end{table}

\vspace{1em}

Figure 3 visualises the CLIP product space: hello's listings sit in the
natural/specialty cluster on the left of the embedding plane while
Colgate's cluster occupies the opposite region, making the low implied
diversion between the merging parties directly visible.

\begin{figure}
\centering
\includegraphics{thesis_body_files/figure-latex/figure1-1.pdf}
\caption{Product space from CLIP embeddings. Each point is one ASIN
positioned by its first two principal components of the joint image+text
embedding. Cluster centroids are labelled C1--C5. hello (red diamonds,
C5) and Colgate (dark blue triangles, C2) occupy different clusters,
implying a lower diversion ratio and a smaller predicted merger effect
under CLIP nesting relative to plain logit.}
\end{figure}

\begin{center}\rule{0.5\linewidth}{0.5pt}\end{center}

\chapter{Empirical Framework}\label{empirical-framework}

\section{Berry (1994) Linearisation}\label{berry-1994-linearisation}

All three models are estimated via the Berry (1994) nested-logit
linearisation:

\[\log\frac{s_{jt}}{s_{0t}} - \rho\,\log s_{j|g,t} = \alpha\,p_{jt} + \boldsymbol{\beta}'\mathbf{x}_{jt} + \text{FE}_t + \xi_{jt}\]

where \(s_{j|g,t} = s_{jt}/\!\sum_{k\in g}s_{kt}\) is the within-nest
share, \(\rho\in[0,1)\) is the nesting intensity, and \(\text{FE}_t\)
are quarter fixed effects (absorbed by within-quarter demeaning). The
underlying choice model is the generalised-extreme-value family of
McFadden (1974, 1978), of which the nested logit is the leading
closed-form member; the error-components representation that motivates
\(\rho\) is due to Cardell (1997), and Train (2009, chs.~3--4) gives a
textbook treatment.

\section{Price Coefficient
Calibration}\label{price-coefficient-calibration}

Valid cost-side instruments for price are unavailable in this
high-frequency online retail dataset (see Section 5). Following the
recommendation of Berry (1994) and standard practice in applied
antitrust simulation, \(\alpha\) is calibrated to match a target median
own-price elasticity:

\[\alpha = \frac{\varepsilon_{\text{target}}}{\operatorname{median}(p_{jt}(1-s_{jt}))} = \frac{-2.62}{8.4502} = -0.310\]

The target \(\varepsilon_{\text{target}} = -2.62\) corresponds to the
brand-level CPG meta-analysis mean of Bijmolt, van Heerde \& Pieters
(2005, p.~141), consistent with structural BLP estimates for analogous
packaged-goods categories (Nevo 2001; Draganska \& Jain 2006).

\textbf{Sequential calibration note.} Since \(\alpha\) is calibrated at
\(\rho=0\), the target elasticity \(-2.62\) is the plain-logit median.
At the profiled nesting parameters, the nested-logit elasticity formula
amplifies own-price sensitivity by roughly \(1/(1-\rho^{*})\), so the
same \(\alpha\) implies median own-price elasticities of \(-3.83\) under
M2 and \(-7.31\) under M3 (computed observation by observation; the
\(1/(1-\rho^{*})\) back-of-envelope gives \(-4.0\) and \(-7.9\)). These
larger magnitudes are not in tension with the calibration target. The
Bijmolt, van Heerde, and Pieters (2005) mean is dominated by brand- and
aggregate-level estimates, whereas the products here are individual
ASINs, which face much closer substitutes than a brand aggregate does --
finer product definitions mechanically imply more elastic demand.
Structural estimates at the product-line level are correspondingly
larger: Draganska and Jain (2006) report own-price elasticities between
\(-2.45\) and \(-6.25\), a range that contains the calibration target
and most of the nested models' implied values. Appendix C reports a
robustness exercise that rescales \(\alpha\) within each nested model so
that the model's own implied median equals \(-2.62\) at its profiled
\(\rho^{*}\): doing so shrinks \(|\alpha|\), inflates Bertrand markups,
and drives the negative-cost rate far above even the plain-logit level.
In other words, ASIN-level demand must be substantially more elastic
than the aggregate benchmark for observed prices to be rationalizable at
all -- the cost side independently validates the more elastic implied
elasticities.

\section{\texorpdfstring{\(\rho^*\) Selection --- Concentrated OLS
Profile}{\textbackslash rho\^{}* Selection --- Concentrated OLS Profile}}\label{rho-selection-concentrated-ols-profile}

For each \(\rho\) on a fine grid \(\{0, 0.01, \ldots, 0.90\}\), the
adjusted dependent variable is formed and \(\boldsymbol{\beta}(\rho)\)
is estimated by OLS. The profiled \(\rho^*\) minimises the residual sum
of squares:

\[\tilde{y}_{jt}(\rho) = \log\frac{s_{jt}}{s_{0t}} - \alpha p_{jt} - \rho\log s_{j|g,t}, \qquad \rho^* = \arg\min_\rho \left\|\tilde{y}(\rho) - \mathbf{X}\hat{\boldsymbol{\beta}}(\rho)\right\|^2\]

\section{Marginal Costs and Merger
Simulation}\label{marginal-costs-and-merger-simulation}

Marginal costs are recovered from the pre-merger Bertrand first-order
condition (Berry 1994):

\[\mathbf{mc} = \mathbf{p} + (\Omega_{\text{pre}} \circ \Delta)^{-1}\mathbf{s}\]

where \(\Omega[j,k] = \mathbf{1}[\text{firm}_j = \text{firm}_k]\) and
\(\Delta\) is the demand Jacobian with elements:

\[\Delta_{jj} = \alpha s_j\!\left[\tfrac{1}{1-\rho} - \tfrac{\rho}{1-\rho}s_{j|g} - s_j\right], \quad \Delta_{jk}^{\text{same}} = \alpha s_j\!\left[-\tfrac{\rho}{1-\rho}s_{k|g} - s_k\right], \quad \Delta_{jk}^{\text{diff}} = -\alpha s_j s_k\]

\textbf{Remark (brand nests and ownership).} Marginal-cost recovery
interacts with the nest assignment in one important, exact way. When
every nest lies wholly inside one firm -- as with brand nests, where
each firm owns complete brands -- the nested-logit Bertrand first-order
conditions admit the same solution as plain logit: markups are uniform
within each firm at \(1/(|\alpha|(1-S_f))\), where \(S_f\) is the firm's
inside share, for \emph{any} value of \(\rho\) (proof in Appendix D).
Consequently, M1 and M2 recover \emph{identical} marginal costs at every
observation, identical Lerner indices, and identical negative-cost
rates. Brand nesting is therefore incapable of repairing the cost side
of the model in this market by construction; only a nest assignment that
cuts across ownership boundaries -- such as the CLIP clusters, which
group ASINs of different firms -- can change recovered costs. The
counterfactual equilibria of M1 and M2 still differ slightly, because
the share \emph{functions} \(s(p)\) differ even though markups at
observed shares coincide.

Marginal costs are averaged per ASIN across pre-merger quarters and held
fixed in both counterfactuals. Holding costs at pre-merger values serves
two purposes: it embodies the standard no-efficiencies assumption of
prospective merger review (Farrell and Shapiro 1990), and it avoids a
mechanical tautology -- if costs were re-inverted in post-merger
quarters, the no-merger Nash equilibrium would recover observed prices
by construction, since the first-order condition is the identity that
defined those costs.

Two Nash equilibria are then computed for every post-merger quarter via
a damped fixed-point iteration on the Bertrand first-order conditions
(Morrow and Skerlos 2011):

\[\mathbf{p}_{n+1} = \lambda\left[\mathbf{mc} - (\Omega \circ \Delta(\mathbf{p}_n))^{-1}\mathbf{s}(\mathbf{p}_n)\right] + (1-\lambda)\,\mathbf{p}_n\]

with \(\lambda = 0.4\), convergence tolerance \(10^{-8}\), and up to
2,000 iterations. The no-merger counterfactual uses pre-merger ownership
\(\Omega_{\text{pre}}\); the merger simulation replaces it with
\(\Omega_{\text{post}}\) (hello ASINs assigned to Colgate's firm
identifier from 2020Q1). The merger price effect is the difference
between the two equilibria,
\(\Delta p = p^{\text{merger}} - p^{\text{no-merger}}\), so that
ownership structure is the only thing that varies between them.

\textbf{Algorithm validation.} The R implementation mirrors the
counterfactual equilibrium routines in PyBLP (Conlon \& Gortmaker 2020)
and is validated against PyBLP's \texttt{compute\_prices()} directly:
for both M1 (plain logit) and M3 (CLIP-nested, \(\rho^*=0.67\)), the
maximum absolute price difference across 847 post-merger observations is
below \(2 \times 10^{-8}\), confirming algorithmic equivalence.

\begin{center}\rule{0.5\linewidth}{0.5pt}\end{center}

\chapter{Results}\label{results}

\section{Demand Estimates}\label{demand-estimates}

Table 4 presents the demand parameters for all three models, separating
the calibrated \(\alpha\) from the estimated \(\boldsymbol{\beta}\) and
the profiled \(\rho^{*}\). Brand nesting yields \(\rho^{*}=0.35\) with a
3.1\% RSS improvement over plain logit; CLIP nesting yields
\(\rho^{*}=0.67\) with an 8.1\% improvement. The stronger profiled
nesting parameter indicates that within-cluster substitution is
substantially more intense when clusters are defined in embedding space
than when they are defined by brand identity.

\begin{table}[H]
\centering
\caption{\label{tab:table2}Demand Parameter Estimates --- Three Models}
\centering
\fontsize{9}{11}\selectfont
\begin{threeparttable}
\begin{tabular}[t]{lccc}
\toprule
Variable & M1: Logit ($\rho$=0) & M2: Brand-Nested ($\rho^*$=0.35) & M3: CLIP-Nested ($\rho^*$=0.67)\\
\midrule
\cellcolor[HTML]{EEF2FF}{\textbf{$\alpha$ (price)}} & \cellcolor[HTML]{EEF2FF}{\textbf{$-$0.310 (a)}} & \cellcolor[HTML]{EEF2FF}{\textbf{$-$0.310 (a)}} & \cellcolor[HTML]{EEF2FF}{\textbf{$-$0.310 (a)}}\\
Package size (g) & 0.0037*** (0.0003) & 0.0037*** (0.0003) & 0.0034*** (0.0003)\\
CLIP $PC_1$ & -0.4034*** (0.0424) & -0.2869*** (0.0418) & -0.3280*** (0.0407)\\
CLIP $PC_2$ & 0.2033*** (0.0374) & 0.2681*** (0.0368) & 0.2103*** (0.0359)\\
CLIP $PC_3$ & -0.1667*** (0.0386) & -0.1479*** (0.0380) & -0.2854*** (0.0370)\\
\addlinespace
CLIP $PC_4$ & 0.2115*** (0.0399) & 0.1574*** (0.0393) & 0.2707*** (0.0382)\\
CLIP $PC_5$ & -0.1652*** (0.0390) & -0.1243*** (0.0384) & -0.1745*** (0.0374)\\
Other brand & 0.0849 (0.1327) & 0.4590*** (0.1307) & 0.2641** (0.1272)\\
Import x freight shock & 0.6767*** (0.2342) & 0.6709*** (0.2306) & 0.7461*** (0.2245)\\
$\rho^*$ & 0.000 (b) & 0.350 (b) & 0.670 (b)\\
\midrule
\addlinespace
Quarter FEs & Yes & Yes & Yes\\
Median implied $\varepsilon_{jj}$ & $-$2.62 & $-$3.83 & $-$7.31\\
RSS drop vs $\rho$=0 & --- & 3.1\% & 8.1\%\\
\bottomrule
\end{tabular}
\begin{tablenotes}
\item \textit{Notes: } 
\item *** p<0.01, ** p<0.05, * p<0.10. FEs by quarter. N=1,164. (a) calibrated alpha; (b) grid-searched rho*. Sections 4.2-4.3.
\end{tablenotes}
\end{threeparttable}
\end{table}

\vspace{1em}

\section{Own-Price Elasticities by
Brand}\label{own-price-elasticities-by-brand}

Table 5 reports median own-price elasticities by brand group for each
model. Two patterns are worth noting. First, within each model,
elasticities scale with price levels -- the nested-logit own-price
elasticity is proportional to \(\alpha p_j\), so premium brands such as
APAGARD and SENSODYNE PRONAMEL are mechanically the most elastic.
Second, and more informative, moving from M1 to M3 amplifies all
elasticities by roughly \(1/(1-\rho^{*})\), but the amplification
interacts with cluster composition: brands sharing a crowded embedding
cluster face the strongest within-nest competition. The merging parties
illustrate the point -- hello's median elasticity rises from \(-1.42\)
under plain logit to \(-4.05\) under CLIP nesting, reflecting the
intense substitution among the natural/specialty products in its cluster
rather than substitution toward Colgate.

\begin{table}[H]
\centering
\caption{\label{tab:table_brand_elas}Median Own-Price Elasticities by Brand Group}
\centering
\fontsize{9}{11}\selectfont
\begin{threeparttable}
\begin{tabular}[t]{lrrrr}
\toprule
Brand & N obs. & M1: Logit & M2: Brand-Nested & M3: CLIP-Nested\\
\midrule
APAGARD & 44 & -4.92 & -6.62 & -14.33\\
Sensodyne & 99 & -3.00 & -4.27 & -8.68\\
SENSODYNE PRONAMEL & 90 & -3.04 & -4.42 & -8.67\\
\textbf{Colgate} & \textbf{224} & \textbf{-2.84} & \textbf{-4.27} & \textbf{-8.24}\\
Tom's of Maine & 85 & -3.03 & -4.38 & -8.07\\
\addlinespace
Other & 220 & -2.70 & -4.09 & -7.98\\
Crest & 219 & -2.11 & -3.18 & -6.17\\
Parodontax & 51 & -1.79 & -2.54 & -5.29\\
Arm \& Hammer & 21 & -1.53 & -1.62 & -4.22\\
\textbf{hello} & \textbf{50} & \textbf{-1.42} & \textbf{-1.88} & \textbf{-4.05}\\
\addlinespace
JASON & 29 & -1.30 & -1.48 & -3.77\\
Orajel & 32 & -0.86 & -1.16 & -2.53\\
\bottomrule
\end{tabular}
\begin{tablenotes}
\item \textit{Notes: } 
\item Median own-price elasticity across all ASIN x quarter observations of each brand group, full estimation panel (1,164 observations). Elasticities use the nested-logit formula of Section 4.4 with the calibrated alpha (-0.310) and each model's profiled rho (0 for M1, 0.35 for M2, 0.67 for M3). Rows sorted by M3 elasticity; bold = merger parties.
\end{tablenotes}
\end{threeparttable}
\end{table}

\vspace{1em}

\section{Diversion Ratios}\label{diversion-ratios}

The merger question turns on where hello's demand goes when its prices
rise. Table 6 reports diversion ratios
\(DR_{j \to k} = -(\partial s_k / \partial p_j) / (\partial s_j / \partial p_j)\)
from hello ASINs, aggregated into fixed reporting buckets
(share-weighted across hello products and pre-merger quarters).
Diversion ratios are the central input of modern unilateral-effects
screening (Farrell and Shapiro 2010; Conlon and Mortimer 2021), and they
enter the U.S. Horizontal Merger Guidelines directly through the
pricing-pressure analysis of differentiated products (U.S. Department of
Justice and Federal Trade Commission 2010, sec.~6.1).

\begin{table}[H]
\centering
\caption{\label{tab:tab_diversion}Diversion from hello by Destination (percent of lost demand)}
\centering
\fontsize{9}{11}\selectfont
\begin{threeparttable}
\begin{tabular}[t]{lrrrr}
\toprule
Model & To Colgate & To rest of hello's CLIP cluster & To other inside goods & To outside good\\
\midrule
M1: Logit & 1.16 & 1.49 & 3.88 & 93.47\\
M2: Brand-Nested & 0.93 & 1.20 & 22.41 & 75.45\\
\cellcolor[HTML]{D6EAD6}{\textbf{M3: CLIP-Nested}} & \cellcolor[HTML]{D6EAD6}{\textbf{0.41}} & \cellcolor[HTML]{D6EAD6}{\textbf{58.14}} & \cellcolor[HTML]{D6EAD6}{\textbf{8.39}} & \cellcolor[HTML]{D6EAD6}{\textbf{33.06}}\\
\bottomrule
\end{tabular}
\begin{tablenotes}
\item \textit{Notes: } 
\item Share-weighted average across hello ASINs and pre-merger quarters (2018Q1--2019Q4), computed from each model's demand Jacobian at observed shares with the calibrated alpha and the model's profiled rho*. Reporting buckets are fixed across models: Colgate-brand products, remaining products in hello's CLIP cluster (C5), all other inside goods, and the outside good. Rows may not sum to 100 due to rounding.
\end{tablenotes}
\end{threeparttable}
\end{table}

The three demand systems tell three different stories about the same
market. Under plain logit, diversion is proportional to market shares,
so 93\% of hello's marginal demand flows to the outside good and only
1.2\% reaches Colgate. Brand nesting redirects substitution toward
hello's own sister ASINs but leaves cross-brand diversion
share-proportional. Under CLIP nesting, 58\% of hello's lost demand
diverts to the other natural/specialty products in its cluster, while
diversion to Colgate falls to 0.4\% -- the low-diversion configuration
that drives the small simulated merger effect in the next subsection.

\section{Merger Counterfactuals}\label{merger-counterfactuals}

Table 7 reports the merger counterfactuals: observed prices, recovered
marginal costs, the no-merger Nash equilibrium, the merger Nash
equilibrium, and the implied price effect \(\Delta p\) for the merging
parties and the largest rivals. Under plain logit, where diversion is
proportional to market shares, the simulated effect for hello is
+1.23\%; brand nesting barely changes it (+1.28\%) because hello's ASINs
sit in their own brand nest and diversion to Colgate is still
share-proportional across nests. Under CLIP nesting, hello (cluster C5)
and Colgate (cluster C2) are distant in product space, diversion between
the parties is small, and the simulated hello effect falls to +0.40\%,
with Colgate's own effect economically negligible.

\begin{table}[H]
\centering
\caption{\label{tab:table3}Merger Counterfactuals: hello-to-Colgate Acquisition (2020Q1)}
\centering
\fontsize{8}{10}\selectfont
\begin{threeparttable}
\begin{tabular}[t]{lrrrrrr}
\toprule
Brand & Obs. $p$ & MC & Nash (pre) & Nash (post) & $\Delta p$ & $\Delta p$ \%\\
\midrule
\addlinespace[0.3em]
\multicolumn{7}{l}{\textit{\textbf{M1: Logit ($\rho$=0)}}}\\
\textbf{\hspace{1em}hello} & \textbf{6.43} & \textbf{3.20} & \textbf{7.10} & \textbf{7.16} & \textbf{0.064} & \textbf{1.233\%}\\
\textbf{\hspace{1em}Colgate} & \textbf{8.45} & \textbf{5.16} & \textbf{8.29} & \textbf{8.30} & \textbf{0.007} & \textbf{0.104\%}\\
\hspace{1em}Tom's of Maine & 10.47 & 7.17 & 10.70 & 10.71 & 0.006 & 0.069\%\\
\hspace{1em}Crest & 8.87 & 5.60 & 9.18 & 9.18 & 0.000 & \vphantom{1} 0.000\%\\
\hspace{1em}Sensodyne & 10.45 & 7.17 & 10.05 & 10.05 & 0.000 & 0.000\%\\
\addlinespace[0.3em]
\multicolumn{7}{l}{\textit{\textbf{M2: Brand-Nested ($\rho^*$=0.35)}}}\\
\textbf{\hspace{1em}hello} & \textbf{6.43} & \textbf{3.20} & \textbf{7.10} & \textbf{7.17} & \textbf{0.065} & \textbf{1.275\%}\\
\textbf{\hspace{1em}Colgate} & \textbf{8.45} & \textbf{5.16} & \textbf{8.30} & \textbf{8.30} & \textbf{0.006} & \textbf{0.104\%}\\
\hspace{1em}Tom's of Maine & 10.47 & 7.17 & 10.71 & 10.71 & 0.006 & 0.069\%\\
\hspace{1em}Crest & 8.87 & 5.60 & 9.18 & 9.18 & 0.000 & 0.000\%\\
\hspace{1em}Sensodyne & 10.45 & 7.17 & 10.07 & 10.07 & 0.000 & 0.000\%\\
\addlinespace[0.3em]
\multicolumn{7}{l}{\textit{\textbf{M3: CLIP-Nested ($\rho^*$=0.67)}}}\\
\textbf{\hspace{1em}hello} & \textbf{6.43} & \textbf{5.20} & \textbf{7.05} & \textbf{7.07} & \textbf{0.020} & \textbf{0.404\%}\\
\textbf{\hspace{1em}Colgate} & \textbf{8.45} & \textbf{5.88} & \textbf{8.59} & \textbf{8.59} & \textbf{0.002} & \textbf{0.018\%}\\
\hspace{1em}Tom's of Maine & 10.47 & 7.19 & 10.72 & 10.72 & 0.002 & 0.021\%\\
\hspace{1em}Crest & 8.87 & 5.63 & 9.18 & 9.18 & 0.000 & 0.000\%\\
\hspace{1em}Sensodyne & 10.45 & 7.28 & 10.29 & 10.29 & 0.000 & 0.000\%\\
\bottomrule
\end{tabular}
\begin{tablenotes}
\item \textit{Notes: } 
\item All prices USD; post-merger averages (2020Q1--2022Q3). Bold = merger parties. Nash (pre): Bertrand-Nash under pre-merger ownership. Nash (post): equilibrium under merged ownership. Dp = Nash (post) minus Nash (pre). MC inverted from pre-merger FOC and held fixed.
\end{tablenotes}
\end{threeparttable}
\end{table}

Figure 4 traces the post-merger window quarter by quarter, showing the
merger Nash equilibrium diverging from the no-merger counterfactual for
hello while the two paths for Colgate remain indistinguishable.

\begin{figure}
\centering
\includegraphics{thesis_body_files/figure-latex/figure2-1.pdf}
\caption{Merger simulation price trajectories (Model 3, CLIP-Nested,
\(\rho^*=0.67\)). Top panel: Nash equilibrium prices under merged and
pre-merger ownership. Observed prices are omitted from the top panel
because their quarter-to-quarter variation (\$4--\$10) would compress
the counterfactual divergence. Bottom panel: merger price effect in
percent,
\((p^{\text{merger}} - p^{\text{no-merger}})/p^{\text{no-merger}} \times 100\).
hello (red) shows a positive and time-varying effect; Colgate (blue)
remains near zero.}
\end{figure}

\vspace{1em}

\section{Cross-Model Comparison}\label{cross-model-comparison}

The central economic-plausibility comparison across the three models is
summarised below. The share of observations with \textbf{negative
implied marginal costs} is the key diagnostic: when the demand system
understates substitutability, the Bertrand inversion attributes
implausibly large markups, pushing recovered costs below zero.
CLIP-based nesting reduces the negative-MC rate from 8.6\% to 1.7\% ---
a five-fold improvement --- while simultaneously delivering the largest
RSS improvement. That M1 and M2 display \emph{exactly} the same
negative-cost rate, cost levels, and Lerner indices is not a
coincidence: because brand nests coincide with ownership boundaries,
brand-nested Bertrand markups are identical to plain-logit markups for
any \(\rho\) (Section 4.4 and Appendix D). The cost side of the model
can only improve when nests cut across firms, which is precisely what
CLIP clustering delivers. The final column translates the price effects
into consumer terms via the nested-logit inclusive-value formula (Small
and Rosen 1981); implied consumer harm under CLIP nesting is roughly a
third of the logit-implied harm.

\begin{table}[H]
\centering
\caption{\label{tab:table4}Cross-Model Comparison: Demand Fit and Economic Plausibility}
\centering
\fontsize{9}{11}\selectfont
\begin{threeparttable}
\begin{tabular}[t]{lcccccccc}
\toprule
Model & $\rho^*$ & Med. $\varepsilon_{jj}$ & RSS drop & Neg. MC \% & Med. Lerner & hello $\Delta p$ & Colgate $\Delta p$ & $\Delta$CS\\
\midrule
M1: Logit & 0.00 & -2.62 & --- & 8.6\% & 38.5\% & +1.23\% & +0.104\% & -25.4\\
M2: Brand-Nested & 0.35 & -3.83 & 3.1\% & 8.6\% & 38.5\% & +1.27\% & +0.104\% & -26.6\\
\cellcolor[HTML]{D6EAD6}{\textbf{M3: CLIP-Nested}} & \cellcolor[HTML]{D6EAD6}{\textbf{0.67}} & \cellcolor[HTML]{D6EAD6}{\textbf{-7.31}} & \cellcolor[HTML]{D6EAD6}{\textbf{8.1\%}} & \cellcolor[HTML]{D6EAD6}{\textbf{1.7\%}} & \cellcolor[HTML]{D6EAD6}{\textbf{28.9\%}} & \cellcolor[HTML]{D6EAD6}{\textbf{+0.40\%}} & \cellcolor[HTML]{D6EAD6}{\textbf{+0.018\%}} & \cellcolor[HTML]{D6EAD6}{\textbf{-8.5}}\\
\bottomrule
\end{tabular}
\begin{tablenotes}
\item \textit{Notes: } 
\item All statistics over the full estimation panel (1,164 ASIN x quarter observations) except merger effects (post-merger quarters 2020Q1--2022Q3). Med. elasticity = median own-price elasticity. RSS drop = reduction in residual sum of squares vs rho=0 from the concentrated OLS profile. Neg. MC = share of observations with negative implied marginal cost from the Bertrand inversion (lower = more economically plausible). Med. Lerner = median (p - mc)/p. hello / Colgate Dp = average post-merger Nash price effect. dCS = mean change in consumer surplus, merger vs no-merger Nash, in cents per 1,000 consumers per quarter (nested-logit log-sum formula). Bold/green row = preferred specification.
\end{tablenotes}
\end{threeparttable}
\end{table}

\begin{center}\rule{0.5\linewidth}{0.5pt}\end{center}

\chapter{Discussion and Limitations}\label{discussion-and-limitations}

The findings should be read with the design's limitations in view; each
is a deliberate trade-off rather than an oversight, and each points to a
concrete extension.

\textbf{Calibrated price coefficient.} Valid cost-side instruments are
unavailable in this high-frequency online retail panel, so \(\alpha\) is
calibrated to an external meta-analytic elasticity rather than
estimated. The choice makes the price-sensitivity assumption
transparent, but it also means \(\alpha\) is common across products and
the level of all markups inherits the calibration. Because \(\alpha\) is
calibrated at \(\rho = 0\), the nested models imply more elastic demand
at their profiled \(\rho^{*}\); the cross-model \emph{comparison} of
fit, negative-cost rates, and merger effects is unaffected, since all
three models share the same \(\alpha\). A natural extension is to
construct differentiation instruments (Gandhi and Houde 2019) from the
CLIP embeddings themselves -- rival products' embedding distances are
plausibly excluded from a product's own demand shock -- and estimate
\(\alpha\) rather than calibrate it.

\textbf{Demand-side cost recovery.} Marginal costs are recovered solely
from the demand system and the Bertrand first-order conditions; without
a supply-side cost equation, true cost variation cannot be separated
from demand shocks absorbed into the inversion. The negative-cost rate
is used here precisely as a \emph{diagnostic} of demand
misspecification, which is robust to this concern, but the recovered
cost levels should not be interpreted structurally.

\textbf{Fixed marginal costs and no efficiencies.} The counterfactuals
hold costs at pre-merger levels, isolating the market-power effect of
the ownership change. If the acquisition generated real marginal-cost
efficiencies, the simulated price effects are upper bounds for the
merged firm's products.

\textbf{Static embeddings and cluster count.} Embeddings are extracted
once per product, which is appropriate for listings that change little,
but would miss repositioning. The number of clusters is selected at
\(K=5\); Appendix A shows the qualitative conclusions -- high
\(\rho^{*}\), low negative-cost rates, small hello effects -- are stable
across \(K = 2\) to \(10\), though point estimates move.

\textbf{Sparse purchase counts.} The crowdsourced panel records 2,787
toothpaste purchases across 1,164 ASIN x quarter cells -- a median of
one purchasing household per cell -- so individual shares carry sampling
noise. The three-quarter presence filter, quarterly aggregation, and the
focus on cross-model \emph{comparisons} (which hold the data fixed)
mitigate this, but individual cell-level quantities should be read with
appropriate caution.

\textbf{External validity.} The analysis covers one category on one
platform around one merger. The pipeline, however, requires only listing
content, purchase counts, and prices -- data available for essentially
any e-commerce category -- which is what makes the approach portable to
other merger settings.

\begin{center}\rule{0.5\linewidth}{0.5pt}\end{center}

\chapter{Conclusion}\label{conclusion}

This paper provides evidence that substitution patterns constructed from
unstructured product content outperform those built from brand labels in
a workhorse merger-simulation framework. Replacing brand nests with
CLIP-embedding clusters in a Berry (1994) nested logit raises demand
fit, cuts the rate of economically implausible negative implied marginal
costs from 8.6\% to 1.7\%, and changes the antitrust answer: the
simulated price effect of the hello--Colgate acquisition for hello falls
from +1.23\% under plain logit to +0.40\% under CLIP nesting, because
the embeddings reveal that the merging parties occupy different regions
of product space and divert little demand to one another.

The broader message for practice is that the inputs to credible merger
simulation need not be limited to structured characteristics data.
Product images and listing text -- observable for any online market at
essentially zero cost -- carry the differentiation consumers act on, and
a pre-trained encoder plus standard clustering suffices to put that
information into the tools antitrust economists already use. The entire
analysis, from embeddings to Nash counterfactuals, is reproducible from
a single document and a public repository, which is intended to make the
method auditable in exactly the settings where it would be applied.

\begin{center}\rule{0.5\linewidth}{0.5pt}\end{center}

\chapter*{Data and Code Availability}\label{data-and-code-availability}
\addcontentsline{toc}{chapter}{Data and Code Availability}

All data, code, and the R Markdown source of this paper are openly
available at
\url{https://github.com/brkuzn/clip-amazon-toothpaste-market}. Knitting
the source document re-runs the entire pipeline -- the joint PCA, the
\(\rho\) grid searches, and the Bertrand--Nash equilibria -- and
regenerates every table and figure from the included data. The
repository also contains the CLIP extraction notebook, the joint-PCA
construction script, and reference copies of all outputs; reproduction
of the shipped joint principal components is byte-exact. Raw
household-level purchase records with demographics are not redistributed
and are available from the Open E-commerce 1.0 repository (Berke et
al.~2024).

\chapter*{References}\label{references}
\addcontentsline{toc}{chapter}{References}

Berke, A., Calacci, D., Mahari, R., Yabe, T., Larson, K., \& Pentland,
A. (2024). Open e-commerce 1.0, five years of crowdsourced U.S. Amazon
purchase histories with user demographics. \emph{Scientific Data}, 11,
491. \url{https://doi.org/10.1038/s41597-024-03329-6}

Berry, S.T. (1994). Estimating discrete-choice models of product
differentiation. \emph{RAND Journal of Economics}, 25(2), 242--262.

Berry, S.T., Levinsohn, J., \& Pakes, A. (1995). Automobile prices in
market equilibrium. \emph{Econometrica}, 63(4), 841--890.

Bijmolt, T.H.A., van Heerde, H.J., \& Pieters, R.G.M. (2005). New
empirical generalizations on the determinants of price elasticity.
\emph{Journal of Marketing Research}, 42(2), 141--156.

Cardell, N.S. (1997). Variance components structures for the
extreme-value and logistic distributions with application to models of
heterogeneity. \emph{Econometric Theory}, 13(2), 185--213.

Compiani, G., Morozov, I., \& Seiler, S. (2026). Demand estimation with
text and image data. \emph{The RAND Journal of Economics}.
\url{https://doi.org/10.1111/1756-2171.70052}

Conlon, C., \& Gortmaker, J. (2020). Best practices for differentiated
products demand estimation with PyBLP. \emph{RAND Journal of Economics},
51(4), 1108--1161.

Conlon, C., \& Mortimer, J.H. (2021). Empirical properties of diversion
ratios. \emph{RAND Journal of Economics}, 52(4), 693--726.

Dosovitskiy, A., Beyer, L., Kolesnikov, A., Weissenborn, D., Zhai, X.,
Unterthiner, T., \ldots \& Houlsby, N. (2021). An image is worth 16x16
words: Transformers for image recognition at scale. \emph{International
Conference on Learning Representations (ICLR)}.

Draganska, M., \& Jain, D.C. (2006). Consumer preferences and
product-line pricing strategies: An empirical analysis. \emph{Marketing
Science}, 25(2), 164--174.

Farrell, J., \& Shapiro, C. (1990). Horizontal mergers: An equilibrium
analysis. \emph{American Economic Review}, 80(1), 107--126.

Farrell, J., \& Shapiro, C. (2010). Antitrust evaluation of horizontal
mergers: An economic alternative to market definition. \emph{The B.E.
Journal of Theoretical Economics}, 10(1), Article 9.

Gandhi, A., \& Houde, J.-F. (2019). Measuring substitution patterns in
differentiated-products industries. NBER Working Paper 26375.

Grigolon, L., \& Verboven, F. (2014). Nested logit or random
coefficients logit? A comparison of alternative discrete choice models
of product differentiation. \emph{Review of Economics and Statistics},
96(5), 916--935.

Hausman, J.A. (1996). Valuation of new goods under perfect and imperfect
competition. In T.F. Bresnahan \& R.J. Gordon (Eds.), \emph{The
Economics of New Goods} (pp.~209--237). University of Chicago Press.

Hotelling, H. (1933). Analysis of a complex of statistical variables
into principal components. \emph{Journal of Educational Psychology},
24(6), 417--441.

Lloyd, S.P. (1982). Least squares quantization in PCM. \emph{IEEE
Transactions on Information Theory}, 28(2), 129--137.

McFadden, D. (1974). Conditional logit analysis of qualitative choice
behavior. In P. Zarembka (Ed.), \emph{Frontiers in Econometrics}
(pp.~105--142). Academic Press.

McFadden, D. (1978). Modelling the choice of residential location. In A.
Karlqvist, L. Lundqvist, F. Snickars, \& J.W. Weibull (Eds.),
\emph{Spatial Interaction Theory and Planning Models} (pp.~75--96).
North-Holland.

Morrow, W.R., \& Skerlos, S.J. (2011). Fixed-point approaches to
computing Bertrand-Nash equilibrium prices under mixed-logit demand.
\emph{Operations Research}, 59(2), 328--345.

Nevo, A. (2000a). A practitioner's guide to estimation of
random-coefficients logit models of demand. \emph{Journal of Economics
\& Management Strategy}, 9(4), 513--548.

Nevo, A. (2000b). Mergers with differentiated products: The case of the
ready-to-eat cereal industry. \emph{RAND Journal of Economics}, 31(3),
395--421.

Nevo, A. (2001). Measuring market power in the ready-to-eat cereal
industry. \emph{Econometrica}, 69(2), 307--342.

Radford, A., Kim, J.W., Hallacy, C., Ramesh, A., Goh, G., Agarwal, S.,
\ldots \& Sutskever, I. (2021). Learning transferable visual models from
natural language supervision. \emph{Proceedings of the 38th
International Conference on Machine Learning (ICML)}, 8748--8763.

Rousseeuw, P.J. (1987). Silhouettes: A graphical aid to the
interpretation and validation of cluster analysis. \emph{Journal of
Computational and Applied Mathematics}, 20, 53--65.

Small, K.A., \& Rosen, H.S. (1981). Applied welfare economics with
discrete choice models. \emph{Econometrica}, 49(1), 105--130.

Train, K.E. (2009). \emph{Discrete Choice Methods with Simulation} (2nd
ed.). Cambridge University Press.

U.S. Department of Justice \& Federal Trade Commission (2010).
\emph{Horizontal Merger Guidelines}. Washington, DC. Issued August 19,
2010.

Werden, G.J., \& Froeb, L.M. (1994). The effects of mergers in
differentiated products industries: Logit demand and merger policy.
\emph{Journal of Law, Economics, \& Organization}, 10(2), 407--426.

\clearpage

\chapter*{Appendix}\label{appendix}
\addcontentsline{toc}{chapter}{Appendix}

\section*{Appendix A: K Sensitivity
Analysis}\label{appendix-a-k-sensitivity-analysis}
\addcontentsline{toc}{section}{Appendix A: K Sensitivity Analysis}

\begin{table}[H]
\centering
\caption{\label{tab:appendix_k_sensitivity}K Sensitivity Analysis: CLIP Cluster Count K = 2--10}
\centering
\fontsize{8}{10}\selectfont
\begin{tabular}[t]{crrcrrrr}
\toprule
$K$ & WCSS & Silhouette & $\rho^*$ & RSS Drop & Neg. MC\% & hello $\Delta p$\% & Colgate $\Delta p$\%\\
\midrule
2 & 1042.0 & 0.199 & 0.60 & 6.1\% & 0.5\% & 1.71\% & 0.015\%\\
3 & 836.7 & 0.256 & 0.66 & 7.8\% & 0.9\% & 3.49\% & 0.013\%\\
4 & 641.2 & 0.330 & 0.57 & 5.7\% & 3.0\% & 3.96\% & 0.041\%\\
\cellcolor[HTML]{D6EAD6}{\textbf{5}} & \cellcolor[HTML]{D6EAD6}{\textbf{456.5}} & \cellcolor[HTML]{D6EAD6}{\textbf{0.388}} & \cellcolor[HTML]{D6EAD6}{\textbf{0.67}} & \cellcolor[HTML]{D6EAD6}{\textbf{8.1\%}} & \cellcolor[HTML]{D6EAD6}{\textbf{1.7\%}} & \cellcolor[HTML]{D6EAD6}{\textbf{0.40\%}} & \cellcolor[HTML]{D6EAD6}{\textbf{0.018\%}}\\
6 & 365.6 & 0.387 & 0.60 & 6.7\% & 3.4\% & 0.59\% & 0.044\%\\
\addlinespace
7 & 322.0 & 0.368 & 0.48 & 4.6\% & 4.4\% & 0.69\% & 0.055\%\\
8 & 279.3 & 0.400 & 0.37 & 2.9\% & 4.7\% & 0.82\% & 0.067\%\\
9 & 247.7 & 0.403 & 0.35 & 2.4\% & 4.7\% & 1.26\% & 0.177\%\\
10 & 230.7 & 0.373 & 0.40 & 3.8\% & 4.6\% & 0.79\% & 0.065\%\\
\bottomrule
\end{tabular}
\end{table}

\emph{Notes:} K-means on 5 CLIP PCs (seed 42, 50 restarts). WCSS =
within-cluster sum of squares. Silhouette = average silhouette width
(higher = better separation). \(\rho^*\) = argmin RSS on grid {[}0,
0.01, \ldots, 0.90{]}. RSS drop = \(100 \times (RSS_0 - RSS^*)/RSS_0\).
Neg. MC = share of ASINs with negative implied marginal cost. hello /
Colgate \(\Delta p\) = average post-merger price effect in percent
(2020Q1--2022Q3). \textbf{Bold/green row = selected specification
(K=5).}

\vspace{1em}

\section*{Appendix B: Model Fit --- RMSE and
MAE}\label{appendix-b-model-fit-rmse-and-mae}
\addcontentsline{toc}{section}{Appendix B: Model Fit --- RMSE and MAE}

\begin{table}[H]
\centering
\caption{\label{tab:appendix_rmse}Model Fit: RMSE / MAE vs. Observed Prices (Post-Merger Quarters 2020Q1--2022Q3)}
\centering
\fontsize{10}{12}\selectfont
\begin{tabular}[t]{lccc}
\toprule
\multicolumn{1}{c}{ } & \multicolumn{3}{c}{RMSE / MAE (USD)} \\
\cmidrule(l{3pt}r{3pt}){2-4}
Brand / Scope & M1: Logit & M2: Brand-Nested & M3: CLIP-Nested\\
\midrule
\textbf{All} & \textbf{1.558 / 0.743} & \textbf{1.556 / 0.746} & \textbf{1.564 / 0.810}\\
hello & 1.579 / 0.836 & 1.580 / 0.837 & 1.533 / 0.799\\
Colgate & 1.144 / 0.559 & 1.142 / 0.560 & 1.149 / 0.652\\
\bottomrule
\end{tabular}
\end{table}

\emph{Notes:} Each cell shows RMSE / MAE in USD against observed
transaction prices. Nash (merger) prices are used as the model
prediction. ``All'' = all 847 ASIN \(\times\) quarter observations in
the post-merger period (2020Q1--2022Q3); hello = 35 obs.; Colgate = 172
obs. Pre-merger marginal costs are held fixed in both Nash simulations.

\vspace{1em}

\section*{Appendix C: Robustness
Checks}\label{appendix-c-robustness-checks}
\addcontentsline{toc}{section}{Appendix C: Robustness Checks}

\begin{table}[H]
\centering
\caption{\label{tab:appendix_robustness}Robustness: Market Size, Calibration Level, and Cost Window}
\centering
\fontsize{8}{10}\selectfont
\begin{tabular}[t]{lrrcrrrr}
\toprule
Scenario & $\lambda$ & $\alpha$ & $\rho^*$ & RSS drop & Neg. MC \% & hello $\Delta p$ & Colgate $\Delta p$\\
\midrule
\textbf{M1 baseline} & \textbf{13.93} & \textbf{-0.310} & \textbf{0.00} & \textbf{0.0\%} & \textbf{8.6\%} & \textbf{+1.233\%} & \textbf{+0.104\%}\\
M1, market size x2 & 27.86 & -0.310 & 0.00 & 0.0\% & 8.4\% & +0.634\% & +0.054\%\\
M1, market size x5 & 69.65 & -0.309 & 0.00 & 0.0\% & 8.4\% & +0.258\% & +0.022\%\\
M1, market size x10 & 139.30 & -0.309 & 0.00 & 0.0\% & 8.4\% & +0.130\% & +0.011\%\\
\textbf{M3 baseline} & \textbf{13.93} & \textbf{-0.310} & \textbf{0.67} & \textbf{8.1\%} & \textbf{1.7\%} & \textbf{+0.404\%} & \textbf{+0.018\%}\\
\addlinespace
M3, market size x2 & 27.86 & -0.310 & 0.67 & 8.1\% & 1.6\% & +0.211\% & +0.010\%\\
M3, market size x5 & 69.65 & -0.309 & 0.67 & 8.1\% & 1.6\% & +0.087\% & +0.004\%\\
M3, market size x10 & 139.30 & -0.309 & 0.67 & 8.1\% & 1.6\% & +0.044\% & +0.002\%\\
M2, alpha recalibrated at rho* & 13.93 & -0.212 & 0.35 & 3.1\% & 20.7\% & +1.810\% & +0.157\%\\
M3, alpha recalibrated at rho* & 13.93 & -0.111 & 0.67 & 8.1\% & 36.3\% & +0.944\% & +0.057\%\\
\addlinespace
M3, MC from 2019Q4 only & 13.93 & -0.310 & 0.67 & 8.1\% & 1.7\% & +0.440\% & +0.022\%\\
\bottomrule
\end{tabular}
\end{table}

\emph{Notes:} Bold rows are the baseline specifications. \textbf{Market
size} rows multiply \(\lambda\) (and hence market size \(M_t\)) by 2, 5,
and 10; \(\alpha\) is recalibrated to the \(-2.62\) target at \(\rho=0\)
within each scenario. Fit and the negative-cost rate are essentially
unchanged; absolute merger effects shrink mechanically as more diversion
flows to the outside good, while the M1-to-M3 ratio of hello effects is
stable. \textbf{Alpha recalibrated at \(\rho^{*}\)} rows hold the
baseline profiled \(\rho^{*}\) fixed and rescale \(\alpha\) so the
model's own implied median elasticity equals \(-2.62\); the resulting
small \(|\alpha|\) inflates Bertrand markups and drives negative-cost
rates far above the plain-logit level, indicating that ASIN-level demand
must be more elastic than the aggregate benchmark to rationalize
observed prices. \textbf{MC from 2019Q4 only} fixes marginal costs at
the last pre-merger quarter instead of the pre-merger average.

\vspace{1em}

\section*{Appendix D: Equivalence of plain-logit and brand-nested
markups}\label{appendix-d-equivalence-of-plain-logit-and-brand-nested-markups}
\addcontentsline{toc}{section}{Appendix D: Equivalence of plain-logit
and brand-nested markups}

\textbf{Proposition.} Consider nested-logit demand with nesting
parameter \(\rho \in [0,1)\) and Bertrand competition, and suppose the
nest partition refines the ownership partition: every nest \(g\) is
wholly owned by a single firm. Then the multiproduct Bertrand markups
solving the first-order conditions at given shares are identical to the
plain-logit (\(\rho = 0\)) markups: for every product \(j\) of firm
\(f\), \(p_j - mc_j = 1/(|\alpha|(1 - S_f))\), where
\(S_f = \sum_{k \in f} s_k\) is the firm's inside share.

\textbf{Proof.} Write \(x_k = p_k - mc_k\). Firm \(f\)'s first-order
condition for product \(j\) is
\(s_j + \sum_{k \in f} x_k\, \partial s_k / \partial p_j = 0\). Under
plain logit, substituting the share derivatives and dividing by
\(\alpha s_j\) gives \(x_j - \sum_{k \in f} s_k x_k = -1/\alpha\), whose
solution is uniform within the firm: \(x_j = \mu_f\) with
\(\mu_f(1 - S_f) = -1/\alpha\). Under nested logit with \(j\) in nest
\(g \subseteq f\), the same substitution gives

\[\frac{x_j}{1-\rho} - \frac{\rho}{1-\rho} \sum_{k \in g} s_{k|g}\, x_k - \sum_{k \in f} s_k x_k = -\frac{1}{\alpha}.\]

Try the uniform candidate \(x_k = \mu_f\) for all \(k \in f\). Because
the entire nest \(g\) belongs to firm \(f\), the within-nest shares sum
to one, \(\sum_{k \in g} s_{k|g} = 1\), so the first two terms collapse
to
\(\mu_f\left(\tfrac{1}{1-\rho} - \tfrac{\rho}{1-\rho}\right) = \mu_f\),
and the condition reduces to \(\mu_f (1 - S_f) = -1/\alpha\) --- the
plain-logit equation. Since the system
\((\Omega \circ \Delta)\,\mathbf{x} = -\mathbf{s}\) is invertible in our
data, this uniform solution is the unique solution, and it does not
depend on \(\rho\).\hfill\(\square\)

The empirical counterpart is exact: across all 1,164 observations, M1
and M2 marginal costs agree to machine precision. The equivalence breaks
whenever a nest spans products of different owners --- the CLIP clusters
do, which is why M3 recovers different (and far more plausible) costs.



\end{document}
